\section{Simulation and Evaluation Setup}
\label{sec:simulation_setup}

\subsection{Simulated sensing system}

The evaluation uses a controlled simulation of a cold-region sensing system motivated by
Antarctic automatic weather stations (AWSs). The final experiment uses six
logical channels, consisting of one mandatory meteorological core channel and five
optional channels. \Cref{tab:sensor_specs} lists the channels and their normalized
scheduling costs. A platform rendering is retained in Supplementary Fig. S2 as
physical motivation, while the channel table defines
the logical-channel configuration evaluated in the experiments.

The simulation keeps the meteorological core channel active, and at most one optional
channel can be active. Core weather measurements remain available while the
optional-channel capacity is allocated among logical channels whose
forecast value changes across event types.

\begin{table}[htbp]
  \centering
  \caption{Logical channels and normalized scheduling costs.}
  \label{tab:sensor_specs}
  \small
  \setlength{\tabcolsep}{3.5pt}
  \begin{tabular}{@{}c >{\raggedright\arraybackslash}p{0.27\linewidth} >{\raggedright\arraybackslash}p{0.33\linewidth} c c c@{}}
    \toprule
    Idx & Logical channel & Main observed variables
        & $p_i$ & $p_i^{\mathrm{peak}}$ & Warm-up \\
    \midrule
    1 & Meteorological core channel
      & wind, air temperature, relative humidity, pressure
      & 0.25 & 0.30 & 0 \\
    2 & Basic radiometer channel
      & solar radiation
      & 0.08 & 0.10 & 0 \\
    3 & Shielded temperature--humidity channel
      & air temperature, relative humidity
      & 0.18 & 0.22 & 0 \\
    4 & Infrared snow-surface-temperature channel
      & snow surface temperature
      & 0.49 & 0.62 & 1 \\
    5 & Laser disdrometer channel
      & particle diameter, particle velocity, simulated particle latent variable
      & 0.49 & 0.62 & 2 \\
    6 & FC4 snow mass-flux channel
      & snow mass flux
      & 0.49 & 0.62 & 1 \\
    \bottomrule
  \end{tabular}
  \vspace{0.5ex}

  \noindent\footnotesize
  The meteorological core channel is mandatory. At each time step, the scheduler may
  activate at most one optional channel and may leave all optional channels
  inactive. The infrared snow-surface-temperature channel, laser disdrometer channel,
  and FC4 snow mass-flux channel are
  event-specific channels for thermal, particle, and flux event types,
  respectively; the basic radiometer channel and shielded temperature--humidity channel
  are additional optional channels. Condition labels are simulator-derived metadata
  used only for training, calibration/validation, or offline evaluation; they are not
  direct channel measurements.
\end{table}


\subsection{Simulated time-series generation}
\label{sec:simulator_construction}

Each seed generates a simulated ground-truth time series with wind, temperature,
humidity, pressure, radiation, snow-surface temperature, particle statistics, and
snow mass flux. The simulation model uses reference distributions from the Antarctic
AWS network and blowing-snow statistics. The meteorological variables are parameterized from the AntAWS
compilation \citep{AntAWS2023}. Event and particle ranges draw on numerical
modeling of atmosphere--snow interaction, Antarctic drifting-snow observations,
and cold-region particle measurements \citep{Amory2020,Sharma2023,Monrad2026}.
Blowing-snow events are produced by a semi-Markov process coupled to wind
exceedance.

The simulation contains particle, flux, and thermal event types. Each event type
changes a different part of the simulated physical state and is observed most
reliably by a different event-specific channel. The laser disdrometer is used for
particle events, the FC4 snow mass-flux channel for flux events, and the infrared
snow-surface-temperature channel for thermal events. Event balance, targets weighted by
event type, and the optional-channel capacity make all three event types matter
in the forecast loss. One fixed optional channel can be effective on average,
but no single optional channel covers the full
forecast objective.

The scheduler input also contains simulated warning signals for the three event
types. These noisy pre-onset signals are assumed available at decision time;
their construction and distinction from condition labels are specified in
Supplementary Section S5. They are simulator-generated and do not constitute a
validated field warning system.

Eight prespecified validation checks assess the simulated series, covering
meteorological structure and blowing-snow behavior \citep{Aloni2024}. They include wind
autocorrelation, event frequency, agreement between simulated and observed
distributions of scalar variables, event durations,
the flux--wind relationship, particle-size--wind correlation, and the wind-speed
power spectrum. All eight criteria pass their prespecified acceptance rules; the
full table and simulation-validation figure are
reported in \ref{apx:generator_validation}. These simulation validation checks document the
simulation construction and remain separate from test policy selection.

\subsection{Metrics, seeds, and statistical summaries}

The primary evaluation uses eight event-balanced evaluation windows of 512 time steps each.
A prespecified rule balances event types and emphasizes
transport periods. The rule
uses event coverage and transport intensity from the simulated ground-truth time series,
not any policy loss, and is applied identically to the PD-PPO scheduler and every baseline. Mean
forecast loss averages \Cref{eq:step_forecast_loss} over the resulting 4096
time steps. The two co-primary endpoints are mean forecast loss and the
macro-averaged normalized forecast loss in \Cref{eq:macro_loss}.
Normalization factors are computed from candidate fixed channel subsets on the
calibration/validation partition, so the loss is
specific to this simulation model. The primary forecaster architecture
and training setup are described in \Cref{sec:framework_protocol}.

Continuous full-partition evaluation examines performance beyond the event-balanced
windows. The same final policy checkpoints and the fixed channel subsets previously
selected on the calibration/validation partition start fresh continuous rollouts and are evaluated at every
test step with a complete eight-step future target. This gives 5,242 scoreable test time steps per seed, covering
$[64750,69992)$; only the last eight
time steps are excluded because their future targets are incomplete. The policies, fixed channel subsets, and normalization factors fitted on the calibration/validation partition are
unchanged, and no schedule is reselected from
this evaluation.

The primary analysis reports 22 post-selection evaluation seeds. Pilot/model-selection seeds 117 and 118 support
a prespecified two-variant architecture comparison using their selected
test-window metrics. The selected configuration is then held fixed for the 22
post-selection evaluation seeds (119--140).
For each seed, the same simulated ground-truth time series is used for the
PD-PPO scheduler, the fixed-schedule baseline,
AoI-priority, round-robin, and random schedules, and all reported baselines. The
validation-selected fixed-schedule baseline is listed before test
evaluation and then uses a constant feasible channel subset without fallback
rotations.

The two pilot/model-selection seeds were used for a prespecified comparison between the
reported warning-conditioned policy and a variant with an additional deterministic
feature encoder. The former was selected before the remaining seeds were run.
Across the 22 post-selection evaluation seeds (119--140), the complete PD-PPO training
configuration achieved lower macro-averaged normalized forecast loss than the
validation-selected fixed-schedule baseline in every seed (mean paired difference 0.0838;
95\% percentile bootstrap confidence interval for the mean paired difference
[0.0708, 0.0967]). The separate 24-seed descriptive aggregate retains the pilot
seeds. The simulation parameters, target weights, baseline
definitions, and reported training settings were held fixed for this expansion.

The main protocol hyperparameters are listed in Supplementary Table S1.
Fixed channel subsets and calibration/validation losses for each seed are listed in
Supplementary Table S9.

The alternative-forecaster analysis uses a ridge forecaster fitted on the same
forecaster-training partition for each seed. The ridge forecaster selects its own
fixed-schedule baseline using calibration/validation windows and then evaluates the unchanged
PD-PPO and fixed-baseline observation sequences. Within this analysis, no policy,
action sequence, or forecaster is updated from its test results. The current
operating-condition label $c_t$ is used offline to select the same
predeclared target weights for every baseline and groups the completed losses for the
macro-averaged normalized forecast loss.

Seed statistical summaries use paired differences because PD-PPO and each baseline
are evaluated on the same simulated ground-truth time series. For a
seed $s$, the paired difference is comparator loss minus PD-PPO loss; a positive
paired difference favors the PD-PPO scheduler. Alongside mean paired differences
and counts of seeds with lower loss, the results report medians and 95\%
percentile bootstrap confidence intervals for the mean paired difference over
seeds.

\subsection{Action-sequence analysis}

The analysis uses the concatenated executed-action sequences from the eight
event-balanced evaluation windows. It reports executed-action entropy,
consecutive-action-pair entropy, condition--executed-action mutual information,
mean normalized Hamming distance between adjacent executed masks, and optional-channel
activation frequency by operating condition. A sequence is classified as fixed-like when
one executed activation mask accounts for at least 95\% of actions and as simple-cycle-like
when the top three masks account for at least 85\% of actions and the best periodic match
up to lag 64 reaches 90\%. The complete behavior-complexity gate also evaluates the
unique executed-action count, both entropy diagnostics, and the predeclared
condition-dependence criteria. These diagnostics include the seven artificial boundaries
between independently reset windows. They support interpretation of the learned schedule; the primary performance
claim remains the paired forecast loss comparison.
